% Figure 14.2: old-fashioned perturbation-theory diagrams for the low-energy
% electron energy shift. Source: physical PDF p. 609 (printed p. 586).
\begingroup
\renewcommand{\thefigure}{14.\arabic{figure}}
\setcounter{figure}{1}
\begin{figure}[tbp]
\centering
\begin{tikzpicture}[
  x=1cm,y=1cm,
  electron/.style={line width=0.72pt},
  photon/.style={line width=0.65pt,decorate,
    decoration={snake,amplitude=1.05pt,segment length=5.0pt}},
  flow/.style={-{Stealth[length=5pt,width=3.8pt]},line width=0.72pt},
  cut/.style={densely dashed,line width=0.65pt}
]
  % First intermediate state: an electron and a photon. Time runs upward.
  \begin{scope}[xshift=-3.15cm]
    \draw[electron] (0,-0.25) -- (0,4.25);
    \draw[photon] (0,0.65) .. controls (1.02,1.05) and
      (1.02,3.35) .. (0,3.70);
    \draw[flow] (0,-0.05) -- (0,0.45);
    \draw[flow] (0,1.25) -- (0,1.92);
    \draw[flow] (0,3.60) -- (0,4.27);
    \draw[cut] (-1.72,2.20) -- (1.72,2.20);
  \end{scope}

  % Second intermediate state: a positron, photon, and initial/final electron.
  \begin{scope}[xshift=3.15cm]
    \coordinate (L) at (-1.25,-0.15);
    \coordinate (R) at (1.30,3.45);
    \draw[electron] (-1.25,-0.15) -- (-1.25,4.35);
    \draw[electron] (1.30,-1.35) -- (1.30,3.45);
    \draw[electron] (L) -- (R);
    \draw[photon] (L) .. controls (0.10,0.25) and
      (1.03,2.10) .. (R);
    \draw[flow] (-1.25,0.50) -- (-1.25,1.18);
    \draw[flow] (1.30,-0.28) -- (1.30,0.40);
    \draw[flow] (0.52,2.30) -- (-0.02,1.55);
    \draw[cut] (-2.15,2.20) -- (2.10,2.20);
  \end{scope}
\end{tikzpicture}
\caption{Old-fashioned perturbation theory diagrams for the low-energy part
of the electron energy shift. Here solid lines are electrons; wavy lines are
photons; the dashed line cuts through particle lines corresponding to the
intermediate states appearing in the first two terms in \eqref{eq:14.3.28}}
\label{fig:14.2}
\end{figure}
\endgroup
